\documentclass{article}
\usepackage{amsthm}

\usepackage[utf8]{inputenc}     % pdfLaTeX
\usepackage[T1]{fontenc}
\usepackage[magyar]{babel}

\title{Óceánia}
\author{Haydu Lőrinc}
\date{Legutóbb frissítve: 2026. Majus 29.}

\newtheorem{definition}{Definíció}

\begin{document}

\maketitle
\tableofcontents
\newpage

\section{Melanézia}
\subsection{Új-Zéland}
\begin{itemize}
    \item 2 nagy + 2 kicsi sziget
    \item Északi szigetek: vulkáni mai napig aktív
    \item Déli szigetek "Déli alpok", jég formálta hegységek
    \item Mérsékelt övezet
    \item Gyűrűk urát itt forgatták
\end{itemize}
\subsection{Új-Guiena}
\begin{itemize}
    \item 2. legnagyobb sziget Grönland után
    \item Összefüggő esőerdő van a területén, majdnem a teljes szigeten
    \item Forró övezet
\end{itemize}
    2. legnagyobb sziget
\section{Mikronézia}
    \begin{itemize}
        \item Mariana-árok \begin{itemize}
            \item 
        \end{itemize}
    \end{itemize}
\begin{itemize}
    \item Atoll: A vulkáni szigetekre az évek alatt korall maradványok rakódnak le. A külső erők lepusztítják a vulkáni hegyet és a helyén ott marad ez a korallgyűrű a, a korall atoll.
\end{itemize}
\section{Polinézia}
\begin{itemize}
    \item Hawaii-szigetek
    \item Aktív és kialudt vulkánok
\end{itemize}
\end{document}